\lecture{12}{Wed 28 Sep 2022 10:30}{}

\begin{theorem}
	One has $A_n \triangleleft S_n$.
\end{theorem}
\begin{proof}
	Let $\sigma \in A_n$. Then for all $\tau\in S_n$ we have to prove $\tau^{-1} \sigma\tau \in A_n$. Thus we need to prove $\tau^{-1} A_n \tau \le A_n$. 

	But if $\sigma \in A_n$, then $\sigma = \omega_1\omega_2\ldots\omega_k$, each $\omega_i$ is a transposition, and $k$ even. Also, $\tau \in S_n$ so $\tau = \nu_1\ldots\nu_l$, with $\nu_i$ transposition.

	\begin{align*}
		\tau^{-1} \sigma\tau
		&= (\nu_1\ldots\nu_l)^{-1} (\omega_1\ldots\omega_k) (\nu_1\ldots\nu_l) \\
		&= \nu_l \ldots\nu_1 \omega_1\ldots\nu_k \nu_1\ldots\nu_l 
	.\end{align*}
	This is a product of transpositions and $2 l + k$ is even, so $\tau^{-1} \sigma\tau \in A_n$. Thus $A_n \triangleleft S_n$.
\end{proof}

To check $|A_n| = \frac{1}{2}n!$ \quad $(n\ge 2)$:
\begin{proof}
	Consider
	\begin{align*}
		\varphi: S_n &\longrightarrow E = \{+1, -1\}  \\
		\sigma &\longmapsto \varphi(\sigma) = (-1)^k
	.\end{align*}
	where $k$ is odd if $\sigma$ is odd, $k$ is even if $\sigma$ is even.

	\begin{claim}
		This is a homomorphism of groups, and is surjective.
	\end{claim}
	If $\sigma, \tau \in S_n$, then $\varphi(\sigma\tau) = (-1)^m .$ where $m$ odd if $\sigma\tau$ is odd, $m$ even if $\sigma\tau$ even.
	
	$\varphi(\sigma)\varphi(\tau) = (-1)^{k+l}$ where $k$ and $l$ depend on parity of $\sigma$ and $\tau$. 
	But $\sigma\tau$ is odd if $\sigma$ and $\tau$ are neither not both even nor not both odd. $\sigma\tau$ is even otherwise. Then $m = k + l \pmod 2$, so $\varphi(\sigma\tau) = \varphi(\sigma)\varphi(\tau)$.

	Then by First Homomorphism Theorem, 
	\[
		S_n / \operatorname{ker} \varphi \cong E
	.\] 
	Also, $\operatorname{ker} \varphi = \{\sigma \in S_n: \varphi(\sigma) = +1\} = \{\sigma \in S_n: \sigma \text{ even}\}  = A_n$.

	Hence $S_n / A_n \cong E$. Hence $|S_n| / |A_n| = |E| = 2$. Hence $|A_n| = \frac{1}{2} |S_n| = \frac{1}{2} n!$.

\end{proof}

\subsection{The group $A_n$ is simple for $n\ge 5$ }

\begin{example}[Simple Groups]
	\begin{enumerate}
		\item $\mathbb{Z}_p$ (cyclic group of prime order)
		\item $A_n \quad (n\ge 5)$
	\end{enumerate}
\end{example}

\begin{theorem}(6.1.2)
	$A_n$ is generated by the $3$-cycles in $S_n$ (when $n \ge 3$ ).
\end{theorem}
\begin{proof}
	If $\sigma \in A_n$, then $A_n$ is a product of an even number of transpositions, say $\sigma = (t_1,t_2)(t_3,t_4)\ldots(t_{2k-1},t_{2k})$, where thet $t_j$ might not be disjoint.

	Consider any pair of these transpositions ( when $k \ge 1$ ). 
	We can relabel $(t_{2i-1}, t_{2i })(t_{2i+1 },t_{2i+2})$ as $(1,2)(3,4)$, or  $(1,2)(1,2)$ or  $(1,2)(1,3)$.
	But  $(1,2)(1,2) = (1,2,3)(1,3,2)$. $(1,2)(3,4) = (1,4,2)(1,4,3)$.

	In all cases these pairs can be written as product of 3-cycles. Because there is a even number of transpositions, we can write $\sigma$ as a product of $3$-cycles.
\end{proof}

\begin{theorem}
	(6.1.8)
	The group $A_5$ is simple.
\end{theorem}
\begin{proof}
	We have $|A_5| = \frac{1}{2}5! = 60$, so all elements $\sigma \in A_5$ have orders dividing $60$ (by Lagrange's Theorem). The possible orders include:
	 \begin{itemize}
		\item $10,15,20,30,60$ : not elements of $S_5$, since cycle decomposition includes $5$-cycle + extra cycles, but no room.
		\item $12$: not element of  $S_5$, since cycle decomposition contain $3$-cycle and $4$-cycle, and not enough room for $7$ elements.
		\item $6$-cycle, by relabeling, looks like $(1,2,3)(4,5)$. This is same as $(1,3)(1,2)(4,5)$ which is odd, so not in $A_5$.
		\item $4$-cycle: by rebelling, looks like $(1,2,3,4) = (1,4)(1,3)(1,2)$ which is odd, so not in $A_5$.
	\end{itemize}
	Then these are elements of order $1,2,3,5$.

	Let's suppose that $N \triangleleft A_5$ and $1< |N| < 60$, and seek a contradiction. Then $N$ contains an element of order $2,3,$ or $5$. We claim that $N$ contains an element of order $3$. Three possibilities:
	\begin{itemize}
		\item  $N$ contains an element of order $3$.
		\item $N$ contains an element of order 2: By relabeling, it can be assumed to be $(1,3)(2,4)$ (note: $(1,2)(1,3)$ has order 3 so is not interesting). But then \[
				(1,2,3,4,5)^{-1}(1,3)(2,4)(1,2,3,4,5) 
				= (1,3)(2,5)
				\in N
		.\] 
		Also, \[
			(1,3)(2,5)(1,3)(2,4)
			= (2, 4, 5) \in N
		.\] 
	\item $N$ contains an element of order $5 $: We can label it as $(1,2,3,4,5)$. Then 
		\[
			\left( (1,3)(2,4) \right) ^{-1} (1,2,3,4,5)\left( (1,3)(2,4) \right) 
			= (1,2,5,3,4)
		.\] 
		Then \[
			(1,2,3,4,5)^{-1}(1,2,5,3,4)
			= (2,4,5)
		.\] 
	\end{itemize}
	So $N$ contains a $3$-cycle, by relabeling we assume to be $(1,2,3)$. We claim that $N$ contains all $3$-cycles. We can replace $(1,2,3)$ by $(1,2,\alpha)$ for any $\alpha \in \{4,5\} $. Now \[
		(2,3,\alpha)^{-1}(1,2,3)^{2}(2,3,\alpha)
		= (1,2,\alpha)
	.\] 
	So whenever $(\alpha,\beta,\gamma) \in N$, then so too does $(\alpha,\beta,\gamma)$ for $\delta \not\in \{\alpha,\beta,\gamma\} $. By switching elements we obtain any $3$-cycle.

	Then $N$ contains all $3$-cycles, and $3$-cycles generate $A_5$, so $N = A_5$, so $A_5$ has no proper normal subgroup. Thus $A_5$ is simple.
\end{proof}
